\section{Conclusão}

Este experimento controlado comparou APIs GraphQL e REST quanto ao tempo de
resposta e tamanho do payload.

\subsection{Respostas às Perguntas de Pesquisa}

\textbf{RQ1 (Tempo de Resposta):} $H_0$ foi rejeitada -- GraphQL foi
significativamente mais rápido que REST (Wilcoxon-Mann-Whitney,
$p < 0,001$), com tempo médio 55\% menor.

\textbf{RQ2 (Tamanho do Payload):} $H_0$ foi rejeitada -- GraphQL retorna
payloads significativamente menores que REST ($p < 0,001$), com redução
de 98,6\% no tamanho (82 bytes contra 6035 bytes).

\subsection{Recomendação}

Com base nos resultados, GraphQL é recomendado para cenários onde:
\begin{itemize}
    \item A eficiência de rede é crítica (aplicações móveis, APIs públicas
          com alto volume de requisições)
    \item O cliente precisa de flexibilidade na seleção de campos
    \item O tempo de resposta é um requisito relevante
\end{itemize}
REST permanece adequado para APIs simples, endpoints estáveis e cenários
onde o overhead de processamento de queries GraphQL não é compensado pela
redução de dados.

\subsection{Trabalhos Futuros}

Sugere-se replicar o experimento com:
\begin{itemize}
    \item Múltiplos objetos experimentais (diferentes APIs)
    \item Consultas de diferentes complexidades (mais campos, \emph{nested queries})
    \item Medição de impacto no servidor (CPU, memória)
    \item Análise com cache habilitado e desabilitado
\end{itemize}
